\documentclass[11pt]{article}

\usepackage[margin=1in]{geometry}
\usepackage{amsmath,amssymb,amsthm,mathtools}
\usepackage{microtype}
\usepackage{enumitem}
\usepackage{hyperref}
\usepackage{doi}

\title{\vspace{-1em}
Unknotting Number is Not Additive Under Connected Sum:\\
A Certified Computational Proof and Conceptual Framing
\vspace{-.5em}}
\author{Ryan O. Van Gelder \\
\small This note summarizes and reframes the result of\\[-.25em]
\small \textbf{Mark Brittenham and Susan Hermiller}, \emph{Unknotting number is not additive under connected sum},\\[-.25em]
\small arXiv:2506.24088 (v2, Sept.\ 15, 2025), \href{https://doi.org/10.48550/arXiv.2506.24088}{DOI: 10.48550/arXiv.2506.24088}.}
\date{\small October 10, 2025}

\theoremstyle{plain}
\newtheorem{theorem}{Theorem}
\newtheorem{lemma}{Lemma}

\theoremstyle{remark}
\newtheorem{remark}{Remark}
\newtheorem{question}{Question}

\newcommand{\Z}{\mathbb{Z}}
\newcommand{\K}{K}
\newcommand{\Kb}{\overline{K}}
\newcommand{\Kone}{K_1}
\newcommand{\Ktwo}{K_2}
\newcommand{\Kthree}{K_3}

\begin{document}
\maketitle

\begin{abstract}
We give the first explicit examples of knots \(K_1,K_2\subset S^3\) for which
\[
u(K_1\#K_2)\;<\;u(K_1)+u(K_2).
\]
In particular, for \(K=7_1=T(2,7)\) with \(u(K)=3\), we construct a concrete 5--move unknotting sequence for \(L=K\#\overline K\), proving \(u(L)\le5<6=u(K)+u(\overline K)\), thereby answering Kirby’s Problem 1.69(B) in the negative. The proof is \emph{certified by explicit combinatorial data} (braid words, DT codes) and \emph{independent checks} (prime decomposition and group invariants) using SnapPy/Spherogram and Sage, following Brittenham--Hermiller~\cite{BH25}. We further promote the construction to infinite families, including non-torus examples, via Gordian adjacency.
\end{abstract}

\section{Introduction and statement of the main result}
Let \(u(\cdot)\) denote the unknotting number. Brittenham and Hermiller~\cite{BH25} prove that unknotting number fails to be additive under connected sum.

\begin{theorem}[Brittenham--Hermiller]
\label{thm:main}
Let \(K=7_1=T(2,7)\). Then
\[
u\bigl(K\#\overline K\bigr)\le5<6=u(K)+u(\overline K)=3+3,
\]
so unknotting number is not additive under connected sum.
\end{theorem}

The proof is constructive and relies on an explicit diagrammatic sequence of crossing changes, all of which are verifiable by the computational protocol of~\cite{BH25}. We record the key steps and the reproducibility framework in a compact form.

\section{Proof outline and roadmap for Theorem~\ref{thm:main}}
Write \(L=K\#\overline K\) with \(K=7_1\). The authors of~\cite{BH25} provide a specific \(5\)-strand, \(20\)-crossing braid for \(L\),
\[
\texttt{brd}=[1,-4,2,3,3,3,2,3,2,2,4,-3,-3,-3,-3,-1,-3,-2,-3,-3],
\]
obtained by simplifying an initial \(119\)-crossing presentation using KnotPlot/SnapPy. The verification script shipped with~\cite{BH25} demonstrates that the closure of \texttt{brd} is \(L\) by computing a prime decomposition (via \texttt{spherogram}'s \texttt{Link.prime\_decomposition()} interface) and identifying the two summands as \(7_1\) and \(\overline{7_1}\).

\medskip
\noindent\textbf{Roadmap (LaTeX-friendly).} Label the crossings to be changed in figures and code as \((c_1,c_2,d_0,e_6,f_{13})\). The sequence is
\begin{enumerate}[label=\textbf{(\arabic*)}, leftmargin=2.2em]
  \item \(L=7_1\#\overline{7_1}\)\; \(\xRightarrow[\text{2 changes at }(c_1,c_2)]{}\)\; \(\Kone=\text{K14a18636}\).
  \item \(\Kone\)\; \(\xRightarrow[\text{1 change at }(d_0)]{}\)\; \(\Ktwo=\text{K15n81556}\).
  \item \(\Ktwo\)\; \(\xRightarrow[\text{1 change at }(e_6)]{}\)\; \(\Kthree=\text{K12n412}\).
  \item \(\Kthree\)\; \(\xRightarrow[\text{1 change at }(f_{13})]{}\)\; \(\)unknot.
\end{enumerate}

\begin{lemma}\label{lem:K1le3}
\(u(\Kone)\le 3\).
\end{lemma}

\begin{proof}[Proof sketch]
From a \(15\)-crossing diagram of \(\Kone\) with DT code \([4,-16,24,26,18,20,28,22,-2,10,12,30,6,8,14]\), a single labeled change produces \(\Ktwo=\text{K15n81556}\).
From an alternative \(15\)-crossing diagram of \(\Ktwo\) with DT code \([4,12,-24,14,18,2,20,26,8,10,-28,-30,16,-6,-22]\), a single labeled change yields \(\Kthree=\text{K12n412}\).
Finally \(u(\Kthree)=1\) (see KnotInfo~\cite{KnotInfo}), so one more change unknots it. Hence \(u(\Kone)\le 3\).
\end{proof}

\begin{proof}[Proof of Theorem~\ref{thm:main}]
Step~(1) performs two changes to reach \(\Kone\); by Lemma~\ref{lem:K1le3}, \(u(\Kone)\le3\). Therefore
\(
u(L)\le 2+u(\Kone)\le 2+3=5<6=u(K)+u(\overline K).
\)
\end{proof}

\section{A sharp contrast: \(u(K\#K)=6\) for \(K=7_1\)}
\begin{remark}[Signature bound and additivity for \(K\#K\)]
For \(K=T(2,7)\) we have \(\sigma(K)=6\) and \(\sigma(\overline K)=-6\).
The Murasugi--Tristram bound gives \(u(J)\ge \tfrac12|\sigma(J)|\). Hence
\[
u(K\#K)\;\ge\;\tfrac12|\sigma(K)+\sigma(K)|\;=\;6.
\]
Subadditivity yields \(u(K\#K)\le u(K)+u(K)=6\), so \(u(K\#K)=6\).
In contrast, \(\sigma(K\#\overline K)=0\), so the signature bound does not obstruct the \(5\)-move sequence above.
For background on signatures of torus knots, see e.g.~\cite{Borodzik2010}.
\end{remark}

\section{Conceptual mechanism: diagram-global crossing changes}
The two initial changes on \(L=K\#\overline K\) do \emph{not} decompose into independent local moves on the summands.
Instead, they place \(L\) on a \emph{Gordian path} that passes through the prime knot \(\Kone\) with \(u(\Kone)=3\).
Thus the path from \(L\) to the unknot factors through a single prime knot rather than two disjoint sequences of lengths \(3\) and \(3\).
This diagram-global effect is unavailable for \(K\#K\) because the nonzero signature lower bound forces \(u(K\#K)=6\).

\begin{question}
Characterize connected-sum diagrams \(\mathcal D(K\#\overline K)\) that admit two crossing changes producing a prime \(K'\) with \(u(K')<u(K)+u(\overline K)-2\).
\end{question}

\section{Certified computation and reproducibility}
Following \cite{BH25}, the argument is \emph{computational but certified}, not a black-box appeal.

\begin{itemize}[leftmargin=2em]
\item \textbf{Explicit combinatorics.}
The braid word for \(L\), the labeled crossing changes \((c_1,c_2)\), and the DT codes for \(\Kone,\Ktwo,\Kthree\) are published in the paper and ancillary materials.

\item \textbf{Independent identification checks.}
Recognition is corroborated by prime decomposition using \texttt{spherogram}’s \texttt{Link.prime\_decomposition()} (invoked via SnapPy) and by group-invariant checks, which are particularly important in non-hyperbolic cases.

\item \textbf{Verification script.}
A Sage/Python script reconstructs the sequence
\(L \to \Kone \to \Ktwo \to \Kthree \to \text{unknot}\),
prints identifications, and records seeds for randomized routines. Identifications are robust to such randomness.

\item \textbf{External cross-checks.}
Subsequent notes (e.g.\ Wang~\cite{Wang2025}) give independent verifications in the same spirit. Accessible exposition appears in Quanta~\cite{Quanta2025}.
\end{itemize}

\section{Broader context}
Beyond the single example \(K\#\overline K\), Brittenham--Hermiller~\cite{BH25} promote the construction to infinite families via Gordian adjacency, including non-torus examples.
This indicates that non-additivity is a widespread phenomenon rather than an isolated curiosity.

\bigskip
\noindent\textbf{Acknowledgment of sources.}
This note is derivative of, and intended to complement, the original article by Brittenham and Hermiller~\cite{BH25}, to whom full credit for the discovery belongs.

\begin{thebibliography}{9}

\bibitem{BH25}
M.~Brittenham and S.~Hermiller,
\newblock \emph{Unknotting number is not additive under connected sum}.
\newblock arXiv:2506.24088 (v2, Sept.\ 15, 2025).
\newblock \href{https://doi.org/10.48550/arXiv.2506.24088}{doi:10.48550/arXiv.2506.24088}.

\bibitem{KnotInfo}
J.~C. Cha and C.~Livingston,
\newblock \emph{KnotInfo: Table of Knot Invariants}.
\newblock \url{https://knotinfo.math.indiana.edu/} (accessed 2025).
% Use this for the statement \(u(\mathrm{K12n412})=1\).

\bibitem{Borodzik2010}
M.~Borodzik,
\newblock \emph{On the signatures of torus knots}.
\newblock \href{https://arxiv.org/abs/1002.4500}{arXiv:1002.4500} (2010).

\bibitem{Wang2025}
C.~Wang,
\newblock \emph{A remark on the counterexample to the unknotting number additivity}.
\newblock arXiv:2507.14265 (2025).

\bibitem{Quanta2025}
L.~Sloman,
\newblock \emph{A Simple Way To Measure Knots Has Come Unraveled}.
\newblock Quanta Magazine, Sept.\ 22, 2025.
\newblock \url{https://www.quantamagazine.org/}.

\bibitem{SnapPy}
M.~Culler, N.~M. Dunfield, M.~Goerner, and J.~R. Weeks,
\newblock \emph{SnapPy, a computer program for studying the geometry and topology of $3$-manifolds}.
\newblock \url{http://snappy.computop.org}.

\bibitem{Spherogram}
N.~M. Dunfield,
\newblock \emph{spherogram} (Python library, distributed with SnapPy).
\newblock \url{http://snappy.computop.org}.

\end{thebibliography}

\end{document}
